\chapter{Cài đặt thực nghiệm và đánh giá}

\section{Mục tiêu đánh giá và khung câu hỏi nghiên cứu}

Mục tiêu của chương này là đánh giá hệ thống trên ba mặt đồng thời: \textit{chất lượng khuyến nghị}, \textit{độ vững vận hành}, và \textit{khả năng giải trình}. Khác với các báo cáo chỉ trình bày kết quả chạy được, chương này đặt trọng tâm vào kiểm định tính hợp lý của lựa chọn kiến trúc dưới ràng buộc MVP.

Bộ câu hỏi nghiên cứu (Research Questions -- RQ) được định nghĩa như sau:

\begin{itemize}
  \item \textbf{RQ1}: Kiến trúc parse-review-publish có giảm lỗi lan truyền vào matching hay không?
  \item \textbf{RQ2}: Thiết kế hard-filter trước semantic scoring có cải thiện tính nhất quán shortlist không?
  \item \textbf{RQ3}: Explanation-by-construction có giúp người dùng nghiệp vụ ra quyết định tốt hơn so với điểm số thuần không?
  \item \textbf{RQ4}: Trong các phương án thay thế phổ biến, cấu hình hiện tại có cân bằng tốt hơn giữa chất lượng, chi phí và bảo trì không?
  \item \textbf{RQ5}: Hệ thống có giữ được hành vi ổn định khi dữ liệu đầu vào suy giảm chất lượng hoặc khi tải tăng đột biến vừa phải không?
\end{itemize}

\section{Thiết kế thực nghiệm end-to-end}

\begin{figure}[H]
\centering
\fbox{\parbox[c][5.1cm][c]{0.9\textwidth}{\centering
TODO: upload image\\
\textbf{Hình 4.1.} Flow thực nghiệm end-to-end: ingest $\rightarrow$ parse/review $\rightarrow$ matching $\rightarrow$ interaction
}}
\caption{Placeholder sơ đồ luồng thực nghiệm end-to-end}
\label{fig:experiment-flow-placeholder}
\end{figure}

Thiết kế thực nghiệm được tổ chức thành các lát cắt tương ứng với vòng đời nghiệp vụ:

\begin{itemize}
  \item \textbf{Slice A -- Data Readiness}: ingest, parse, normalize, review gate.
  \item \textbf{Slice B -- Candidate Generation}: lexical, dense, hybrid merge.
  \item \textbf{Slice C -- Ranking \& Explanation}: hard-filter, scoring, explanation payload.
  \item \textbf{Slice D -- ATS Workflow}: apply, invite, trạng thái và audit trail.
  \item \textbf{Slice E -- Reliability}: hành vi khi lỗi cục bộ xuất hiện trong từng tầng.
\end{itemize}

Mỗi slice được chạy theo ba kịch bản dữ liệu: bình thường, nhiễu vừa, và biên bất lợi. Cách tổ chức này giúp đánh giá hệ thống theo điều kiện gần thực tế hơn so với một kịch bản đơn tuyến tính.

\section{Mapping RQ -> Metric -> Evidence}

\begin{table}[H]
\caption{Ma trận liên kết câu hỏi nghiên cứu với độ đo và bằng chứng}
\centering
\begin{tabularx}{\textwidth}{p{0.10\textwidth}p{0.27\textwidth}p{0.24\textwidth}X}
\toprule
\textbf{RQ} & \textbf{Mục tiêu kiểm định} & \textbf{Metric chính} & \textbf{Bằng chứng quan sát}\
\\
\midrule
RQ1 & Chặn lỗi dữ liệu lan vào matching & Tỷ lệ hồ sơ bị giữ ở review gate; tỷ lệ lỗi downstream & Nhật ký trạng thái parse/review/publish theo lô dữ liệu.\\
RQ2 & Nhất quán shortlist nghiệp vụ & Tỷ lệ cặp vi phạm constraint xuất hiện top-$k$; độ ổn định thứ hạng gần kề & So sánh kết quả trước/sau hard-filter.\\
RQ3 & Hiệu quả giải trình & Tỷ lệ quyết định nghiệp vụ có thể truy nguyên thành phần điểm & Mẫu explanation payload và phản hồi thao tác người dùng.\\
RQ4 & Cân bằng kiến trúc & Chi phí suy luận, độ trễ, độ phức tạp vận hành, chất lượng top-$k$ & Phân tích phản thực theo baseline thay thế.\\
RQ5 & Độ bền vận hành & Tỷ lệ hoàn tất pipeline khi có lỗi cục bộ; thời gian phục hồi tương đối & Log lỗi có cấu trúc và kết quả retry/fallback.\\
\bottomrule
\end{tabularx}
\end{table}

Ma trận trên là cơ sở để tránh đánh giá cảm tính. Mỗi kết luận đều phải truy ngược được về một hoặc nhiều metric cùng bằng chứng tương ứng.

\section{Tiêu chí và độ đo đánh giá}

\subsection{Nhóm độ đo vận hành}

\begin{itemize}
  \item \textbf{Pipeline Completion Rate}: tỷ lệ hồ sơ đi trọn từ ingest đến trạng thái sẵn sàng.
  \item \textbf{Stage Latency Profile}: độ trễ theo tầng (parse, retrieval, ranking, workflow action).
  \item \textbf{State Transition Validity}: tỷ lệ chuyển trạng thái hợp lệ trong ATS workflow.
  \item \textbf{Recovery Effectiveness}: tỷ lệ lỗi cục bộ được phục hồi qua retry/fallback.
\end{itemize}

\subsection{Nhóm độ đo chất lượng khuyến nghị}

\begin{itemize}
  \item \textbf{Constraint Precision}: mức độ loại đúng cặp không khả thi ở tầng hard-filter.
  \item \textbf{Top-$k$ Utility}: mức hữu ích shortlist cho hành động nghiệp vụ tiếp theo.
  \item \textbf{Explanation Utility}: khả năng giải thích hỗ trợ quyết định thay vì gây nhiễu nhận thức.
\end{itemize}

\subsection{Nhóm độ đo định lượng roadmap}

Dù chưa dùng làm tiêu chí chốt hiện tại, bộ chỉ số định lượng chuẩn được chuẩn bị cho giai đoạn mở rộng gồm: precision@$k$, recall@$k$, nDCG@$k$, MRR, và calibration error theo nhóm vị trí. Việc chuẩn bị trước bộ chỉ số này giúp giảm rủi ro thay đổi tiêu chí giữa chừng khi bước vào benchmark có nhãn.

\section{Kết quả quan sát theo từng slice}

\subsection{Slice A: Data Readiness}

Kết quả cho thấy review gate giúp cô lập hồ sơ có rủi ro parser trước khi vào matching công khai. Tác động quan sát được là giảm lỗi lan truyền và giảm số lần phải xử lý phản hồi sai lệch ở downstream. Trong môi trường dữ liệu không đồng nhất, lợi ích này lớn hơn chi phí tăng thêm một bước review.

\subsection{Slice B: Candidate Generation}

Candidate generation theo hướng hybrid cho độ phủ ổn định hơn khi truy vấn có biến thể ngôn ngữ hoặc mô tả kỹ năng không chuẩn. Khi truy vấn mang tín hiệu từ vựng mạnh, lexical giữ ưu thế; khi truy vấn giàu ngữ nghĩa mô tả, dense bù đắp tốt. Sự bổ trợ này tạo nền tốt hơn cho tầng ranking sau đó.

\subsection{Slice C: Ranking \& Explanation}

Thiết kế hard-filter trước scoring giúp loại bỏ mâu thuẫn ``điểm cao nhưng không khả thi''. Ở phần explanation, cấu trúc theo thành phần cho phép người dùng nghiệp vụ truy nguyên nhanh lý do xếp hạng và xác định điểm cần hiệu chỉnh dữ liệu. Đây là khác biệt đáng kể so với các danh sách điểm tổng hợp thiếu giải thích.

\subsection{Slice D: ATS Workflow}

State machine giúp các chuyển trạng thái diễn ra nhất quán, giảm lỗi ``nhảy bước'' và giúp audit dễ hơn. Khi có tranh chấp hoặc cần hậu kiểm, log actor + timestamp + pre/post state cung cấp đường truy vết đầy đủ.

\subsection{Slice E: Reliability under degradation}

Khi mô phỏng lỗi cục bộ ở parse hoặc retrieval, kiến trúc tách lớp cho phép hệ thống suy giảm mềm: một số tác vụ bị chậm hoặc tạm hoãn nhưng không làm toàn bộ pipeline ngừng hoạt động. Đây là thuộc tính quan trọng cho hệ thống tuyển dụng vận hành liên tục.

\section{Phân tích phản thực (counterfactual) theo baseline}

\subsection{Baseline A: Keyword-only pipeline}

Nếu dùng keyword-only làm lõi, hệ thống có thể giảm chi phí suy luận và đơn giản hóa hạ tầng. Tuy nhiên, shortlist dễ lệch về các hồ sơ ``khớp chữ'' thay vì ``khớp nghĩa'', đặc biệt trong các JD mô tả dài hoặc đa ngữ cảnh. Về dài hạn, bias từ cách viết CV/JD sẽ ảnh hưởng chất lượng matching nhiều hơn mức tiết kiệm chi phí ngắn hạn.

\subsection{Baseline B: Dense blackbox end-to-end}

Nếu dùng dense blackbox end-to-end, chất lượng semantic retrieval có thể cải thiện trong một số kịch bản. Tuy nhiên, hai vấn đề vận hành xuất hiện: khó giải trình cho quyết định nghiệp vụ và khó phân rã lỗi khi đầu vào suy giảm chất lượng. Khi chưa có bộ nhãn dày, rủi ro tối ưu sai mục tiêu (offline tốt, online không bền) là đáng kể.

\subsection{Baseline C: KG-heavy + multi-agent}

Nếu đưa KG-heavy và orchestration đa tác tử ngay từ MVP, hệ thống có thể tăng năng lực suy luận tri thức dài hạn. Đổi lại, độ phức tạp triển khai tăng mạnh: ontology design, data governance, đồng bộ tri thức, và kiểm thử liên tác tử. Trong điều kiện MVP, chi phí này khó biện minh so với lợi ích tức thì.

\subsection{Kết luận phản thực}

Các baseline thay thế đều có điểm mạnh cục bộ, nhưng kiến trúc hiện tại đạt cân bằng tốt hơn cho mục tiêu MVP: chất lượng đủ dùng, giải thích được, vận hành ổn định, và có lộ trình nâng cấp từng lớp.

\section{Threats to Validity}

\subsection{Internal Validity}

Rủi ro nội tại lớn nhất nằm ở confounding giữa chất lượng dữ liệu parser và chất lượng ranking. Nếu không tách hai lớp này khi phân tích, kết luận về scorer có thể sai lệch. Cơ chế review gate và logging theo tầng được dùng để giảm rủi ro này, nhưng không loại bỏ hoàn toàn.

\subsection{Construct Validity}

Một số khái niệm như ``độ hữu ích của giải thích'' khó lượng hóa bằng một chỉ số đơn. Vì vậy, construct này được đo bằng tổ hợp chỉ báo (khả năng truy nguyên quyết định, thời gian phản hồi thao tác, mức nhất quán hành động nghiệp vụ). Dù vậy, vẫn còn rủi ro đo lường chưa bao phủ đủ chiều sâu nhận thức người dùng.

\subsection{External Validity}

Kết quả hiện tại phản ánh dữ liệu và quy trình ở phạm vi MVP. Khi mở rộng sang ngành khác hoặc thị trường ngôn ngữ khác, phân bố dữ liệu và kỳ vọng tuyển dụng có thể thay đổi. Do đó, mọi kết luận tổng quát hóa cần đi kèm tái đánh giá theo miền.

\subsection{Conclusion Validity}

Do thiếu tập nhãn quy mô lớn, các kết luận định lượng hiện chưa đạt mức suy diễn thống kê mạnh như các benchmark học thuật chuẩn. Vì vậy, kết luận chương này được đóng khung ở mức ``evidence-informed engineering'' thay vì ``statistically final claims''.

\section{Bảng tổng hợp mức độ đáp ứng hiện tại}

\begin{table}[H]
\caption{Tổng hợp mức đáp ứng theo trục đánh giá MVP}
\centering
\begin{tabularx}{\textwidth}{p{0.25\textwidth}p{0.15\textwidth}X}
\toprule
\textbf{Trục đánh giá} & \textbf{Mức} & \textbf{Diễn giải}\
\\
\midrule
Độ đúng quy trình & Tốt & Checkpoint trạng thái giúp kiểm soát vòng đời dữ liệu/nghiệp vụ nhất quán.\\
Độ phủ candidate & Khá-Tốt & Hybrid generation giảm bỏ sót so với một hướng đơn lẻ.\\
Độ hữu ích shortlist & Khá & Hard-filter + scoring thành phần tạo danh sách gợi ý có thể hành động.\\
Khả năng giải trình & Khá & Explanation payload hỗ trợ truy nguyên, còn dư địa mở rộng attribution sâu.\\
Độ bền vận hành & Khá & Tách lớp giúp suy giảm mềm khi lỗi cục bộ, cần tăng telemetry ở tải cao.\\
\bottomrule
\end{tabularx}
\end{table}

\noindent
\textbf{Nhận xét tổng quan:} hệ thống đạt mục tiêu MVP theo định hướng kỹ thuật thực dụng. Điểm mạnh nằm ở kiến trúc có kiểm soát và khả năng truy vết; điểm cần đầu tư tiếp theo là chuẩn hóa benchmark định lượng có nhãn và tăng chiều sâu đánh giá online.

\section{Roadmap nghiên cứu và nâng cấp ba giai đoạn}

\subsection{Near-term (0--3 tháng)}

\begin{itemize}
  \item Chuẩn hóa pipeline nhãn top-$k$ theo chuyên gia tuyển dụng.
  \item Bổ sung dashboard theo tầng: parse quality, retrieval drift, ranking stability.
  \item Chuẩn hóa bộ câu hỏi phản hồi explanation từ người dùng nghiệp vụ.
\end{itemize}

\subsection{Mid-term (3--9 tháng)}

\begin{itemize}
  \item Triển khai benchmark định lượng: precision@$k$, recall@$k$, nDCG@$k$, MRR.
  \item Chạy ablation cho từng lớp (hard-filter, semantic scoring, rerank signal).
  \item Thí điểm calibration theo nhóm vị trí để giảm sai lệch liên miền.
\end{itemize}

\subsection{Long-term (9+ tháng)}

\begin{itemize}
  \item Thử nghiệm kiểm soát KG augmentation ở domain hẹp.
  \item Đánh giá khả thi multi-agent orchestration cho các tác vụ phân vai rõ (không thay lõi ngay).
  \item Mở rộng đánh giá fairness/ethics theo bối cảnh AI hiring có kiểm chứng dữ liệu thực địa.
\end{itemize}

\begin{figure}[H]
\centering
\fbox{\parbox[c][4.8cm][c]{0.9\textwidth}{\centering
TODO: upload image\\
\textbf{Hình 4.3.} Roadmap nghiên cứu 3 giai đoạn và điều kiện kích hoạt nâng cấp
}}
\caption{Placeholder sơ đồ roadmap nâng cấp theo giai đoạn}
\label{fig:roadmap-3phase-placeholder}
\end{figure}

\section{Protocol thực nghiệm chi tiết cho từng RQ}

Để tái lập đánh giá, mỗi RQ được gắn với protocol thao tác cụ thể thay vì mô tả tổng quát. Protocol gồm bốn phần: dữ liệu đầu vào, thao tác chạy, điều kiện quan sát, và tiêu chí kết luận.

\begin{table}[H]
\caption{Protocol thực nghiệm rút gọn theo RQ}
\centering
\begin{tabularx}{\textwidth}{p{0.10\textwidth}p{0.23\textwidth}p{0.23\textwidth}X}
\toprule
\textbf{RQ} & \textbf{Input profile} & \textbf{Thao tác chính} & \textbf{Tiêu chí kết luận}\\
\midrule
RQ1 & CV/JD có mức nhiễu parser khác nhau & Chạy ingest-parse-review qua nhiều batch & Review gate phải chặn được phần lớn hồ sơ thiếu trường trọng yếu.\\
RQ2 & Bộ truy vấn có và không có hard constraints & Chạy matching với/không có hard-filter & Tỷ lệ cặp vi phạm constraint ở top-$k$ giảm rõ khi bật hard-filter.\\
RQ3 & Tập kết quả có explanation payload đầy đủ & So sánh quyết định thao tác với/không có explanation & Explanation giúp tăng khả năng truy nguyên và nhất quán thao tác.\\
RQ4 & Cùng tập dữ liệu, nhiều baseline kiến trúc & Chạy phân tích phản thực theo giả định kỹ thuật & Cấu hình hiện tại phải chứng minh lợi thế cân bằng, không chỉ một metric đơn lẻ.\\
RQ5 & Kịch bản lỗi cục bộ theo tầng & Inject lỗi có kiểm soát và quan sát phục hồi & Hệ thống phải suy giảm mềm, không sập toàn tuyến.\\
\bottomrule
\end{tabularx}
\end{table}

Protocol chi tiết theo RQ giúp tăng độ tin cậy kết luận và giảm tranh luận cảm tính khi so sánh phương án.

\section{Phân tích độ nhạy (sensitivity analysis) theo tham số kiến trúc}

Một sai lầm thường gặp là xem trọng số scoring hoặc độ sâu rerank như hằng số bất biến. Trên thực tế, các tham số này có độ nhạy cao theo loại vị trí tuyển dụng. Do đó, đánh giá cần xem biến thiên kết quả khi thay đổi có kiểm soát các tham số chính:

\begin{itemize}
  \item Trọng số lexical/dense/rule adjustment.
  \item Độ sâu candidate set trước rerank.
  \item Ngưỡng hard-filter ở các trường điều kiện cứng.
  \item Chính sách normalization score theo batch.
\end{itemize}

\begin{table}[H]
\caption{Khung sensitivity analysis cho các tham số trọng yếu}
\centering
\begin{tabularx}{\textwidth}{p{0.28\textwidth}p{0.20\textwidth}X}
\toprule
\textbf{Tham số} & \textbf{Khoảng thử} & \textbf{Kỳ vọng quan sát}\\
\midrule
$\alpha$ lexical weight & 0.2 -- 0.6 & Tăng $\alpha$ giúp truy vấn giàu từ khóa, nhưng có thể giảm bao phủ ngữ nghĩa.\\
$\beta$ dense weight & 0.3 -- 0.7 & Tăng $\beta$ giúp semantic recall, nhưng cần kiểm soát chi phí và độ ổn định explanation.\\
Rerank depth & 20 -- 200 ứng viên & Độ sâu lớn có thể tăng chất lượng top-$k$ nhưng tăng độ trễ đáng kể.\\
Constraint threshold & mềm -- chặt & Ngưỡng chặt giảm false positive nhưng có thể tăng false rejection ở trường hợp biên.\\
\bottomrule
\end{tabularx}
\end{table}

Phân tích độ nhạy không nhằm tìm một bộ tham số ``đúng tuyệt đối'', mà nhằm xác định vùng tham số an toàn cho vận hành.

\section{Ablation framework cho kiến trúc nhiều tầng}

Ablation là công cụ quan trọng để định lượng đóng góp của từng lớp trong pipeline. Với JobConnect, ablation nên chạy theo thứ tự tăng dần độ phức tạp:

\begin{itemize}
  \item Cấu hình A0: lexical + hard-filter.
  \item Cấu hình A1: A0 + dense retrieval.
  \item Cấu hình A2: A1 + hybrid merge policy.
  \item Cấu hình A3: A2 + rerank signal.
  \item Cấu hình A4: A3 + explanation-by-construction đầy đủ.
\end{itemize}

\begin{figure}[H]
\centering
\fbox{\parbox[c][4.8cm][c]{0.9\textwidth}{\centering
TODO: upload image\\
\textbf{Hình 4.4.} Khung ablation A0-A4 và thứ tự bật thành phần kiến trúc
}}
\caption{Placeholder sơ đồ ablation framework cho pipeline nhiều tầng}
\label{fig:ablation-framework-placeholder}
\end{figure}

Khung ablation này giúp trả lời câu hỏi thực dụng: thành phần nào tạo giá trị thực, thành phần nào chỉ tăng độ phức tạp mà lợi ích biên thấp.

\section{Khuyến nghị chuẩn hóa báo cáo kết quả cho hội đồng}

Để báo cáo có tính học thuật cao và dễ phản biện, phần kết quả nên chuẩn hóa theo mẫu:

\begin{itemize}
  \item Nêu rõ giả thiết kỹ thuật trước mỗi bảng/biểu đồ.
  \item Tách riêng kết quả quan sát và diễn giải nguyên nhân.
  \item Với mỗi kết luận, chỉ ra metric và nguồn evidence tương ứng.
  \item Tránh tuyên bố tổng quát vượt quá phạm vi dữ liệu đánh giá.
\end{itemize}

Chuẩn hóa mẫu báo cáo giúp giữ nhất quán giữa các vòng cập nhật luận văn, đồng thời giảm rủi ro lập luận vòng tròn hoặc suy diễn quá mức từ dữ liệu hạn chế.

\section{Rủi ro triển khai khi mở rộng từ MVP lên production}

Khi chuyển từ MVP sang production, ba rủi ro kỹ thuật xuất hiện sớm:

\begin{itemize}
  \item \textbf{Scale risk}: độ trễ tăng do số lượng hồ sơ và truy vấn tăng cùng lúc.
  \item \textbf{Governance risk}: thay đổi policy tuyển dụng nhanh hơn vòng cập nhật scoring rule.
  \item \textbf{Model lifecycle risk}: phiên bản mô hình/embedding không được quản trị nhất quán.
\end{itemize}

Để giảm rủi ro, đề xuất tối thiểu là: version hóa policy và model, tách môi trường benchmark khỏi môi trường vận hành, và thêm cơ chế canary cho thay đổi scoring logic.

\section{Governance thực nghiệm và khả năng tái lập kết quả}

Một báo cáo học thuật mạnh cần khả năng tái lập. Với hệ thống ứng dụng như JobConnect, tái lập không chỉ là chạy lại code mà còn là tái lập đúng ngữ cảnh dữ liệu, phiên bản policy và cấu hình pipeline.

\begin{table}[H]
\caption{Experiment card tối thiểu cho mỗi vòng đánh giá}
\centering
\begin{tabularx}{\textwidth}{p{0.22\textwidth}X}
\toprule
\textbf{Trường bắt buộc} & \textbf{Nội dung cần ghi nhận}\\
\midrule
Dataset slice & Nguồn CV/JD, khoảng thời gian snapshot, tỷ lệ hồ sơ thiếu trường, mức nhiễu parser.\\
Version bundle & Policy version, model version, index version, cấu hình hard-filter/rerank.\\
RQ target & Câu hỏi nghiên cứu chính của lượt chạy và giả thuyết cần kiểm chứng.\\
Metrics & Bộ metric chính/phụ, ngưỡng chấp nhận, phương pháp tổng hợp.\\
Outcome notes & Kết quả định lượng, quan sát định tính, quyết định giữ/đổi cấu hình.\\
Risk notes & Các bất thường, giới hạn dữ liệu, mức độ tin cậy kết luận.\\
\bottomrule
\end{tabularx}
\end{table}

Chuẩn hóa experiment card giúp hội đồng hoặc nhóm vận hành kiểm tra được đường suy luận từ dữ liệu đến kết luận, giảm rủi ro ``kết luận đúng do may mắn dữ liệu''.

\section{Phân tích chi phí vận hành theo baseline và ngưỡng đầu tư}

Ngoài metric chất lượng, luận văn cần làm rõ câu hỏi: khi nào nên đầu tư sang cấu hình phức tạp hơn? Ta định nghĩa chi phí kỳ vọng theo truy vấn:

\begin{equation}
E[\text{Cost/query}] = c_r + d \cdot c_{rr} + c_o
\end{equation}

trong đó $c_r$ là chi phí retrieval nền, $d$ là rerank depth, $c_{rr}$ là chi phí rerank trên mỗi ứng viên, và $c_o$ là overhead observability/governance.

Lợi ích biên của cấu hình nâng cấp được xem xét qua:

\begin{equation}
\Delta U = \Delta Q - \gamma \cdot \Delta \text{Cost}
\end{equation}

Nếu $\Delta U > 0$ ổn định trên nhiều batch và nhiều RQ, nâng cấp có thể được chấp nhận. Nếu không, cấu hình hiện tại vẫn là lựa chọn tối ưu trong ngữ cảnh MVP.

\begin{table}[H]
\caption{Khung quyết định đầu tư theo lợi ích biên}
\centering
\begin{tabularx}{\textwidth}{p{0.24\textwidth}p{0.22\textwidth}X}
\toprule
\textbf{Tình huống} & \textbf{Dấu hiệu định lượng} & \textbf{Khuyến nghị}\\
\midrule
Nâng cấp rerank sâu & $\Delta Q$ nhỏ, $\Delta$Cost lớn & Giữ cấu hình hiện tại, chỉ tăng depth cục bộ theo nhóm truy vấn khó.\\
Thêm attribution phức tạp & Explainability tăng nhẹ, latency tăng đáng kể & Hoãn triển khai rộng, thử nghiệm trên domain hẹp có nhu cầu audit cao.\\
Tối ưu hybrid fusion & $\Delta Q$ và stability cùng tăng, cost tăng thấp & Ưu tiên triển khai, vì cải thiện utility tổng thể.\\
\bottomrule
\end{tabularx}
\end{table}

\section{Phân tích lỗi theo error bucket và giá trị cho roadmap}

Phân tích lỗi nên đi theo ``error bucket'' thay vì chỉ liệt kê case lẻ:

\begin{itemize}
  \item \textbf{Bucket E1 -- Data quality}: thiếu trường kỹ năng, kinh nghiệm, hoặc metadata ngữ cảnh.
  \item \textbf{Bucket E2 -- Retrieval miss}: candidate set bỏ sót ứng viên phù hợp do lexical/dense lệch pha.
  \item \textbf{Bucket E3 -- Ranking misorder}: candidate đúng có mặt nhưng thứ hạng không ổn định ở top-$k$.
  \item \textbf{Bucket E4 -- Explanation inconsistency}: điểm thành phần và narrative không đồng nhất.
  \item \textbf{Bucket E5 -- Workflow friction}: kết quả đúng kỹ thuật nhưng khó dùng trong thao tác ATS.
\end{itemize}

Mỗi bucket cần ánh xạ trực tiếp vào roadmap:
\begin{itemize}
  \item E1 ưu tiên near-term vì tác động diện rộng.
  \item E2/E3 ưu tiên mid-term với ablation và sensitivity có kiểm soát.
  \item E4/E5 là nền cho long-term governance và trust.
\end{itemize}

\begin{figure}[H]
\centering
\fbox{\parbox[c][4.8cm][c]{0.9\textwidth}{\centering
TODO: upload image\\
\textbf{Hình 4.5.} Bản đồ error bucket và ánh xạ sang roadmap 3 giai đoạn
}}
\caption{Placeholder bản đồ lỗi hệ thống và ưu tiên xử lý theo giai đoạn}
\label{fig:error-bucket-roadmap-placeholder}
\end{figure}

Khung error bucket giúp biến phần ``hạn chế'' từ mô tả thụ động thành chương trình cải tiến chủ động có thứ tự ưu tiên rõ ràng.

\section{Khung báo cáo negative results và quyết định dừng sớm}

Trong nghiên cứu ứng dụng, không phải mọi hướng nâng cấp đều tạo cải thiện đáng kể. Việc ghi nhận ``negative results'' (kết quả không cải thiện) có vai trò quan trọng vì giúp tránh lặp lại thử nghiệm tốn kém và tăng chất lượng quyết định đầu tư.

Đối với JobConnect, một vòng thử nghiệm nên được đánh dấu là negative result khi thỏa đồng thời:
\begin{itemize}
  \item Cải thiện chất lượng không vượt ngưỡng tối thiểu đã định trước.
  \item Độ trễ hoặc chi phí tăng vượt ngưỡng vận hành chấp nhận.
  \item Explainability không tốt hơn hoặc khó truy nguyên hơn cấu hình hiện tại.
\end{itemize}

Khi xuất hiện negative result, hành động phù hợp không phải luôn là ``thử lại ngay'', mà là phân loại nguyên nhân:
\begin{itemize}
  \item \textbf{Do dữ liệu}: cần cải thiện parse/normalize trước khi đánh giá lại mô hình.
  \item \textbf{Do cấu hình}: giữ mô hình nhưng đổi độ sâu rerank hoặc trọng số fusion.
  \item \textbf{Do giả thuyết sai}: đóng nhánh thử nghiệm và chuyển tài nguyên sang hướng khác.
\end{itemize}

\begin{table}[H]
\caption{Quy tắc dừng sớm cho thử nghiệm nâng cấp}
\centering
\begin{tabularx}{\textwidth}{p{0.24\textwidth}p{0.2\textwidth}X}
\toprule
\textbf{Tín hiệu} & \textbf{Mức độ} & \textbf{Quyết định}\\
\midrule
$\Delta Q$ thấp liên tiếp nhiều batch & Cao & Dừng thử nghiệm, quay về baseline ổn định, lưu biên bản negative result.\\
Chi phí tăng nhanh hơn lợi ích & Trung bình-Cao & Giảm phạm vi thử nghiệm hoặc hạ cấu hình để xác định ngưỡng hiệu quả.\\
Kết quả dao động mạnh giữa batch & Trung bình & Tăng cỡ mẫu và chuẩn hóa protocol trước khi kết luận.\\
\bottomrule
\end{tabularx}
\end{table}

Khung dừng sớm giúp Chương 4 không chỉ trả lời ``hướng nào tốt'' mà còn trả lời ``hướng nào nên dừng'', từ đó tăng tính thực dụng và độ tin cậy của roadmap nghiên cứu.

\section{Kết luận chương}

Chương 4 đã chuyển đánh giá từ mức ``hệ thống chạy được'' sang mức ``hệ thống hợp lý về mặt kiến trúc và vận hành dưới ràng buộc MVP''. Phân tích phản thực cho thấy cấu hình hiện tại không cực đại thuật toán, nhưng tối ưu hơn về cân bằng chất lượng--chi phí--giải trình. Phần threats to validity và roadmap ba giai đoạn cung cấp nền tảng rõ ràng cho các vòng nghiên cứu và nâng cấp tiếp theo.
